\documentclass[11pt]{article}
\usepackage[T1]{fontenc}
\usepackage{lmodern}
\usepackage[margin=1in]{geometry}
\usepackage{amsmath,amssymb,amsthm,mathtools}
\usepackage{microtype}
\usepackage[colorlinks=true,linkcolor=blue,citecolor=blue,urlcolor=blue]{hyperref}
\hypersetup{pdftitle={An asymptotic lower bound for the number of irrational even beta values},pdfauthor={twinheterodox},pdfsubject={Irrationality counts for even Dirichlet beta values}}
\newtheorem{theorem}{Theorem}[section]
\newtheorem{proposition}[theorem]{Proposition}
\newtheorem{lemma}[theorem]{Lemma}
\newtheorem{corollary}[theorem]{Corollary}
\theoremstyle{remark}
\newtheorem{remark}[theorem]{Remark}
\newcommand{\Z}{\mathbb Z}
\newcommand{\Q}{\mathbb Q}
\newcommand{\R}{\mathbb R}
\newcommand{\C}{\mathbb C}
\newcommand{\Lcal}{\mathcal L}
\newcommand{\sgn}{\operatorname{sgn}}
\newcommand{\Res}{\operatorname{Res}}
\newcommand{\dd}{\,\mathrm d}
\DeclareMathOperator{\lcm}{lcm}
\DeclareMathOperator{\rank}{rank}
\DeclareMathOperator{\Span}{Span}
\numberwithin{equation}{section}
\title{\Large An asymptotic lower bound for the number of irrational even beta values}
\author{twinheterodox}
\date{November 17, 2026}

\begin{document}
\maketitle

\begin{abstract}
For even $s$, let $I(s)$ count the irrational values among
$\beta(2),\beta(4),\ldots,\beta(s)$. We prove
\[
 \liminf_{\substack{s\to\infty\\2\mid s}}
 \frac{I(s)}{\sqrt{s\log\log s/\log s}}
 \geq 2\sqrt{\frac{\zeta(2)\zeta(3)}{3\zeta(6)}}
 \approx 1.6098018250 .
\]
Here $\beta$ denotes the Dirichlet beta function and $\zeta$ the Riemann
zeta function. The bound improves
the growth rate in our earlier
counting result by a factor $\sqrt{\log\log s}$. It has the same growth
rate as the recent theorem of Fischler, Lai, Sprang and Zudilin for
irrational odd zeta values, with $1/\sqrt3$ of their leading constant.
The conclusion remains valid after any fixed initial set of even beta
values is omitted.
\end{abstract}

\section{The counting problem and the character twist}
The Dirichlet beta function is
\[
 \beta(u)=\sum_{v=0}^{\infty}\frac{(-1)^v}{(2v+1)^u}
         =L(u,\chi_4),\qquad
 \chi_4(j)=
 \begin{cases}
  0,&2\mid j,\\
  1,&j\equiv1\pmod4,\\
 -1,&j\equiv3\pmod4.
 \end{cases}
\]
Its positive odd integer values are nonzero rational multiples of powers
of $\pi$. The even values have a different arithmetic character. Our aim
is to count how many of them must be irrational, without specifying
which ones.

Write
\[
 I(s)=\#\{i:2\leq i\leq s,\ 2\mid i,\ \beta(i)\notin\Q\},\qquad
 \kappa=\frac{\zeta(2)\zeta(3)}{\zeta(6)},\qquad
 Y(s)=\sqrt{\frac{s\log\log s}{\log s}}.
\]
Here and throughout, $s$ denotes the largest allowed index, rather than
the number of entries in the list.

\begin{theorem}\label{thm:main}
For every $\varepsilon>0$, all sufficiently large even integers $s$ satisfy
\[
 I(s)\geq\left(2\sqrt{\frac{\kappa}{3}}-\varepsilon\right)Y(s).
\]
\end{theorem}

Deleting finitely many indices changes $I(s)$ by a bounded amount, whereas
$Y(s)\to\infty$. Thus the same lower limit holds if, for example, the list
starts with $\beta(4)$. The theorem gives neither a numerical threshold for
$s$ nor a lower bound for the dimension of the span of these values.

Our earlier explicit counting theorem \cite[Theorem~1.1]{THcount} gives
\[
 I(s)\geq\left\lfloor\pi\sqrt{\frac{2}{135}}
                      \sqrt{\frac{s}{\log s}}\right\rfloor
 \qquad(s\geq160,\quad 2\mid s).
\]
Theorem~\ref{thm:main} improves the growth rate by a factor
$\sqrt{\log\log s}$. Its starting endpoint depends on $\varepsilon$
and is not computed here, whereas the earlier bound has the stated
explicit range. The endpoint certificates for that range are in the
\texttt{even-beta-counts/checks} directory of the project repository~\cite{beta4repo}.

Fischler, Lai, Sprang and Zudilin \cite[Theorem~1.1]{FLSZ} prove the
corresponding odd-zeta counting theorem with leading constant
$2\sqrt\kappa$. Their proof combines elimination of presumed irrational
values, the rational functions of Lai and Yu \cite{LY}, and a Hankel
determinant argument for signed discrete measures. Kumar and Tosi
\cite[Theorem~1]{KT} state a counting bound of order $\sqrt{s/\log s}$
for the noncritical values of a real Dirichlet character, including
$\chi_4$. Their stated constant at conductor four is
\[
 c(4)=\left(\frac{\kappa}{2}
  \log\frac{\sqrt{4e^2+16}+4}{\sqrt{4e^2+16}-4}\right)^{1/2}.
\]
The proof below does not use that theorem or its denominator estimates.

Our construction is a character twist of \cite[Sections~2--6]{FLSZ}.
On a grid whose denominator is divisible by four, the signs $\chi_4(j)$
combine Hurwitz zeta values into beta values. Only odd $j$ contribute.
We therefore put zeros in the numerator only at the fractions used by
the character and their integer translates. This retains every elimination
weight and reduces the leading logarithmic cost of the numerator by a
factor of three. We can consequently use a grid cutoff larger by a
factor of $\sqrt3$. Together with the density $1/3$ of usable denominators,
this explains the ratio $1/\sqrt3$ between the two leading constants.

An even pole order selects the even indices under reflection. The
combined grid weights are odd, so we also check that the reflected
measure has the opposite residues needed for nonvanishing. The
arithmetic normalization, analytic test polynomials and signed-measure
criterion follow \cite{FLSZ} with the restricted numerator.

The argument assumes familiarity with elementary prime valuations,
polynomial linear algebra and residues of meromorphic functions. We give
the signed-measure criterion in full before applying it to the character
weights.

\section{The available grids}\label{sec:grids}
The grid family is chosen to supply many elimination weights while
keeping the rational function's numerator cost under control. For $B>1$,
let
\begin{align*}
 \Psi_B&=\{b\in\Z_{>0}:\varphi(b)\leq B\},\\
 \mathcal Q_B&=\{q\in\Psi_B:4\mid q\},\\
 \Theta_B&=\left\{\frac aq:q\in\mathcal Q_B,\ 1\leq a<q,\ (a,q)=1\right\},
\end{align*}
where $\varphi$ is Euler's totient function, and put $h_B=|\Theta_B|$.
If $q\in\mathcal Q_B$, $1\leq j<q$ and $\chi_4(j)\ne0$, reducing $j/q$ cancels
only an odd divisor. Its reduced denominator $b$ is still divisible
by four, and $\varphi(b)\leq\varphi(q)$ because $b\mid q$. Thus
$j/q$ and $1-j/q$ both belong to $\Theta_B$. Conversely, each reduced
fraction in $\Theta_B$ occurs with nonzero character weight on its
own grid. For example, on the $q=12$ grid, $3/12$ and its reflection
reduce to $1/4$ and $3/4$, while $2/12$ carries weight zero.

Using these same denominators, put
\[
 A_1(B)=\prod_{q\in\mathcal Q_B}q^{\varphi(q)},\qquad
 A_2(B)=\prod_{q\in\mathcal Q_B}\prod_{p\mid q}p^{\varphi(q)/(p-1)},\qquad
 A(B)=A_1(B)A_2(B).
\]
These are integers since $p-1$ divides $\varphi(q)$ for $p\mid q$.
Our starting point is the inverse-totient estimate of
\cite[Proposition~2.2]{LY}, also recorded in \cite[Lemma~2.1]{FLSZ},
\begin{equation}\label{eq:totient-count}
 |\Psi_B|=(\kappa+o(1))B .
\end{equation}

\begin{lemma}\label{lem:grid-count}
The usable denominators satisfy
\[
 |\mathcal Q_B|=\left(\frac\kappa3+o(1)\right)B
 \qquad(B\to\infty).
\]
\end{lemma}
\begin{proof}
For real $x>0$, let $F(x)$ count all positive integers $b$ with
$\varphi(b)\leq x$, let $g(x)$ count the odd such integers, and let $T(x)$
count those divisible by four. Separating the power of two in $b$ gives
\[
 F(x)=2g(x)+T(x),\qquad
 T(x)=\sum_{j\geq1}g(x/2^j).
\]
The two copies of $g$ correspond to $b$ odd and $b$ twice an odd integer;
in both cases the totient is unchanged. It follows that
\[
 2T(x)=F(x/2)+T(x/2),\qquad
 T(x)=\sum_{j\geq1}2^{-j}F(x/2^j).
\]
These sums are finite for fixed $x$, since $F(y)=0$ for $y<1$.
By \eqref{eq:totient-count}, $F(y)/y\to\kappa$ and $F(y)\leq Cy$ for
$y\geq1$, with an absolute constant $C$. Dividing the last sum by $x$
and applying dominated convergence with majorant $C4^{-j}$ yields
$T(x)/x\to\kappa\sum_{j\geq1}4^{-j}=\kappa/3$.
\end{proof}

The same count determines how many numerator factors we pay for. Each
usable denominator contributes exactly $\varphi(q)$ reduced fractions.
Writing $T(t)$ for the count in Lemma~\ref{lem:grid-count}, partial
summation in the totient variable gives
\begin{equation}\label{eq:grid-size}
 h_B=\sum_{q\in\mathcal Q_B}\varphi(q)
     =B T(B)-\int_0^B T(t)\dd t
     =\left(\frac{\kappa}{6}+o(1)\right)B^2 .
\end{equation}
In particular $h_B+1\leq B^2$ for all sufficiently large $B$, since
$\kappa<6$. Every $q\in\Psi_B$ satisfies $q\ll B\log\log B$
\cite[(2.2)]{LY}. Hence $\log A_1(B)\leq h_B(1+o(1))\log B$.
For the reverse bound discard the denominators with $\varphi(q)<B/\log B$.
By \eqref{eq:grid-size} they contribute $o(B^2)$ to $h_B$, and every
remaining $q$ satisfies $\log q\geq\log B-\log\log B$. Thus
\begin{equation}\label{eq:A1}
 \log A_1(B)=\left(\frac{\kappa}{6}+o(1)\right)B^2\log B .
\end{equation}
The corresponding coefficient for the full product is $\kappa/2$
\cite[Lemma~2.2]{FLSZ}. For $A_2$ the full product is an upper bound, so
\cite[Lemma~2.2]{FLSZ} gives, for all sufficiently large $B$,
\begin{equation}\label{eq:A2}
 \log A_2(B)\leq 10 B^2(\log\log B)^2 .
\end{equation}

\section{The twisted functional}
\label{sec:functional}
Fix a sufficiently large even integer $s$, and take $B$ sufficiently large
with $B^2<s/100$. We shall eventually choose $B$ in terms of $s$.
Until Section~\ref{sec:finish}, these two parameters remain fixed and
$n$ tends to infinity through positive even integers. Write $\Theta=\Theta_B$,
$h=h_B$ and $A=A(B)$.

For $(t)_k=t(t+1)\cdots(t+k-1)$, define
\begin{equation}\label{eq:kernel}
 R_n(t)=A^n(n!)^{s-h}
        \frac{\prod_{\theta\in\Theta}(t+\theta)_n}{(t)_{n+1}^{s}}.
\end{equation}
This uses the kernel of \cite[(2.1)]{FLSZ} with even $s$ and the
restricted numerator grid $\Theta$. It is positive
for $t>0$, has degree $nh-(n+1)s$, and satisfies
\begin{equation}\label{eq:reflection}
 R_n(-n-t)=R_n(t).
\end{equation}
Indeed, reflection pairs $\theta$ with $1-\theta$ in the numerator;
$n$ is even, and the sign from the denominator disappears because $s$
is even.

Suppose, for a contradiction, that $I(s)<|\mathcal Q_B|$. There are then
integers $w_q$, not all zero, such that
\begin{equation}\label{eq:elimination}
 \sum_{q\in\mathcal Q_B}w_q q^i=0
 \quad\text{whenever }2\leq i\leq s,\quad2\mid i,\quad\beta(i)\notin\Q.
\end{equation}
This homogeneous rational system has fewer equations than unknowns; clearing
denominators gives integral weights. Choose a positive integer $A_0$ clearing
the denominators of the remaining, rational, beta values in the same range.
Both $A_0$ and the $w_q$ are fixed as $n$ varies.

Writing $\zeta(i,\theta)=\sum_{v\geq0}(v+\theta)^{-i}$ for the Hurwitz
zeta function, the character identity underlying the construction is,
for $i>1$ and $4\mid q$,
\begin{equation}\label{eq:character}
 \sum_{j=1}^q\chi_4(j)\zeta(i,j/q)=q^i\beta(i).
\end{equation}
For example, the grid $q=4$ gives
$\zeta(i,1/4)-\zeta(i,3/4)=4^i\beta(i)$.
In general, the integers in the Hurwitz sums are $qv+j$, and
$\chi_4(qv+j)=\chi_4(j)$; this proves \eqref{eq:character} by absolutely
convergent summation.

For a polynomial $P$, set
\begin{equation}\label{eq:functional}
 \Lcal_n(P)=A_0\sum_{q\in\mathcal Q_B}w_q\sum_{j=1}^q\chi_4(j)
  \sum_{v=0}^{\infty}
  R_n(v+j/q)
  \left[P(v+j/q+n/2)+P(-v-j/q-n/2)\right].
\end{equation}
We call $P$ admissible if
\begin{equation}\label{eq:admissible}
 \deg P\leq(n+1)s-nh-2.
\end{equation}
This condition makes the summand $O(v^{-2})$ and guarantees absolute
convergence. The proof uses the same moment matrix in two bases. The
first makes its scaled determinant an integer. The second bounds the
absolute value of that integer by a number less than one. To obtain a
contradiction, we allow the size $m$ to vary and prove that at least one
permitted size has nonzero determinant. We write
\begin{equation}\label{eq:det}
 \Delta_{n,m}=\det\bigl(\Lcal_n(t^{i+j})\bigr)_{0\leq i,j<m},
 \qquad 2m\leq(n+1)s-nh.
\end{equation}

\section{Arithmetic normalization}\label{sec:arithmetic}
Let $d_n=\lcm(1,\ldots,n)$. The role of the numerator zeros is more than
analytic: they cancel the possible extra denominators in the finite
Hurwitz tails. We spell out this step because twisting the grids must not
silently change the integerizer.

\subsection{The elementary derivative bounds}
The following facts are the building blocks used in
\cite[Section~3]{FLSZ}. A direct proof also shows why their arithmetic
does not depend on the parity of $s$.

\begin{lemma}\label{lem:blocks}
For a fixed positive integer $n$, each of the following functions $F$
satisfies $d_n^\ell F^{(\ell)}(k)/\ell!\in\Z$ for all $\ell\geq0$
at the indicated integer evaluation points:
\begin{enumerate}
\item $(t+a)_r/r!$, with $a\in\Z$ and $0\leq r\leq n$, at every $k\in\Z$;
\item $b^n\prod_{p\mid b}p^{\lfloor n/(p-1)\rfloor}(t+a/b)_n/n!$,
      with $(a,b)=1$ and $b>1$, at every $k\in\Z$;
\item $n!(t+j)/(t)_{n+1}$, with $0\leq j\leq n$, at $k=-j$.
\end{enumerate}
\end{lemma}
\begin{proof}
The case $r=0$ and Taylor orders above the degree are immediate.
For the first assertion with $r\geq1$, a Taylor coefficient of order $\ell\leq r$
is a sum of products obtained by deleting $\ell$ factors from $r$
consecutive integers, divided by $r!$. For any prime $p$ and integer
$v\geq1$, the original progression has at least $\lfloor r/p^v\rfloor$
factors divisible by $p^v$. Deleting $\ell$ factors loses at most $\ell$
of them. Hence each product has valuation at least
\[
 v_p(r!)-\ell\lfloor\log_p r\rfloor
 \geq v_p(r!)-\ell v_p(d_n).
\]
This proves integrality after multiplication by $d_n^\ell$.

For the second assertion, write its numerator as
$\prod_{i=0}^{n-1}(bt+a+bi)$. If $p\nmid b$, this arithmetic progression
has at least $\lfloor n/p^v\rfloor$ multiples of each $p^v$; the same
deletion argument applies, and the extra factors $b^\ell$ are integral.
If $p\mid b$, the prefactor $p^{\lfloor n/(p-1)\rfloor}$ alone clears
$n!$, since $v_p(n!)\leq\lfloor n/(p-1)\rfloor$. All the numerator
coefficients are integers, proving the assertion at these primes too.

For the last assertion set $t=-j+u$. The constant coefficient is
$(-1)^j\binom nj$, and the remaining factor is
\[
 \prod_{\substack{0\leq k\leq n\\k\ne j}}
     \left(1+\frac{u}{k-j}\right)^{-1}.
\]
Each Taylor coefficient of order $\ell$ is a sum of products of $\ell$
reciprocals of nonzero integers of absolute value at most $n$.
Multiplication by $d_n^\ell$ clears them.
\end{proof}

\subsection{The coefficients and the finite tails}

\begin{lemma}\label{lem:arithmetic}
If an admissible $U\in\Q[t]$ satisfies
\begin{equation}\label{eq:derivative-integrality}
 \frac{d_n^\ell}{\ell!}U^{(\ell)}(k)\in\Z
 \qquad(k\in\Z,\ \ell\geq0),
\end{equation}
then $d_n^s\Lcal_n(U)\in\Z$.
\end{lemma}
\begin{proof}
Write
\[
 \widetilde R_n(t)=R_n(t)\left[U(t+n/2)+U(-t-n/2)\right]
   =\sum_{k=0}^n\sum_{i=1}^s\frac{a_{i,k}}{(t+k)^i}.
\]
The decomposition has no polynomial part by admissibility. The arithmetic
building-block argument of \cite[(3.5)--(3.6)]{FLSZ} gives
\begin{equation}\label{eq:coeff-int}
 d_n^{s-i}a_{i,k}\in\Z\qquad(1\leq i\leq s,\ 0\leq k\leq n).
\end{equation}
For clarity, $(t+k)^s\widetilde R_n(t)$ factors into $s$ copies of
$n!(t+k)/(t)_{n+1}$, the blocks
\[
 b^n\prod_{p\mid b}p^{\lfloor n/(p-1)\rfloor}
       \frac{(t+a/b)_n}{n!}
 \quad(b\in\mathcal Q_B,\ 1\leq a<b,\ (a,b)=1),
\]
and $U(t+n/2)+U(-t-n/2)$, together with an integer factor.
Lemma~\ref{lem:blocks} and the hypothesis on $U$ give
$d_n^\ell F^{(\ell)}(-k)/\ell!\in\Z$ for each factor.
The leftover constant is integral because
$n\varphi(b)/(p-1)\geq\varphi(b)\lfloor n/(p-1)\rfloor$ for every
$p\mid b$. Applying the product rule
to the Taylor coefficient of order $s-i$ proves \eqref{eq:coeff-int}.
Neither this calculation nor its building-block estimates require $s$ to
be odd. The shifts of $U$ are integral because $n$ is even.

Reflection now has the opposite parity to the odd-zeta construction:
\[
 a_{i,n-k}=(-1)^i a_{i,k}.
\]
Consequently $\rho_i=\sum_{k=0}^n a_{i,k}$ vanishes for every odd $i$.
For $0<\theta<1$, summing the partial fractions therefore gives
\begin{equation}\label{eq:hurwitz-expansion}
 \sum_{v\geq0}\widetilde R_n(v+\theta)
 =\rho_{0,\theta}+\sum_{\substack{2\leq i\leq s\\2\mid i}}
       \rho_i\zeta(i,\theta),\qquad
 \rho_{0,\theta}=-\sum_{k=1}^n\sum_{l=0}^{k-1}
                    \sum_{i=1}^s\frac{a_{i,k}}{(l+\theta)^i}.
\end{equation}
The $i=1$ terms are combined in finite sums before taking their limit.
Their common harmonic divergence has coefficient $\rho_1=0$, and the
remaining shifted endpoint terms tend to zero. Thus no value
$\zeta(1,\theta)$ is introduced.

We claim $d_n^s\rho_{0,\theta}\in\Z$ for each $\theta=j/q$ with
$q\in\mathcal Q_B$ and $\chi_4(j)\ne0$. Such $\theta$ and $1-\theta$
belong to $\Theta$. Fix $0\leq l<k\leq n$ and a prime $p$, and put
$e=v_p(d_n)$. If $v_p(l+\theta)\leq e$, each term
$d_n^s a_{i,k}/(l+\theta)^i$ is $p$-integral by \eqref{eq:coeff-int}.
Otherwise, evaluate the partial fractions at
$t_0=-k+l+\theta$. This is a numerator zero: the block
$(t+1-\theta)_n$ vanishes there, since
$1\leq k-l\leq n$. Hence
\[
 \sum_{i=1}^s\frac{d_n^s a_{i,k}}{(l+\theta)^i}
 =-\sum_{\substack{0\leq k'\leq n\\k'\ne k}}\sum_{i=1}^s
       \frac{d_n^s a_{i,k'}}{(l+\theta+k'-k)^i}.
\]
Since $v_p(k'-k)\leq e<v_p(l+\theta)$, the denominator
$(l+\theta+k'-k)^i$ has valuation $i\,v_p(k'-k)\leq ie$.
Equation~\eqref{eq:coeff-int} gives $v_p(d_n^s a_{i,k'})\geq ie$,
so every term on the right is $p$-integral. This proves the claim
prime by prime, after summing the finite tails.

Finally apply \eqref{eq:character} to \eqref{eq:hurwitz-expansion} and
sum with weights $w_q$. The irrational beta terms disappear by
\eqref{eq:elimination}; the rational ones are cleared by $A_0$, while
\eqref{eq:coeff-int} clears their coefficients after multiplication by
$d_n^s$. The finite tails have just been cleared by the same multiplier.
\end{proof}

The command \mbox{\texttt{python3 checks/verify\_count\_examples.py}} reproduces
the finite-tail cancellations for $B=2$, $s=8$ and $n=2,4$. Run it from
\texttt{asymptotic-even-beta-counts/} in the project repository~\cite{beta4repo}.

\begin{proposition}\label{prop:integerizer}
For $i\geq0$, define
\[
 E_{n,i}=(n!)^{\lfloor i/n\rfloor}(i-n\lfloor i/n\rfloor)!,\qquad
 H_{n,m}=\frac{d_n^{sm}}{\prod_{i=0}^{m-1}E_{n,i}^2}.
\]
For every $m$ allowed in \eqref{eq:det}, $H_{n,m}\Delta_{n,m}\in\Z$.
\end{proposition}
\begin{proof}
Use $U_i(t)=(t)_i/E_{n,i}$ as a basis of polynomials of degree below $m$.
It is a product of normalized Pochhammer polynomials of degrees at most
$n$: divide the $i$ factors into blocks of $n$ and a final shorter block.
Lemma~\ref{lem:blocks} and the product rule, as in
\cite[Proposition~3.2]{FLSZ}, give \eqref{eq:derivative-integrality} for
each product $U_iU_j$. Lemma~\ref{lem:arithmetic} shows that all entries of
$(d_n^s\Lcal_n(U_iU_j))$ are integers. The triangular change of basis
has determinant $\prod_{i<m}E_{n,i}^{-1}$. Applying it to the rows and
columns gives
\[
 \det\bigl(d_n^s\Lcal_n(U_iU_j)\bigr)_{i,j<m}
   =d_n^{sm}\left(\prod_{i<m}E_{n,i}^{-1}\right)^2\Delta_{n,m}
   =H_{n,m}\Delta_{n,m}.
\]
\end{proof}

\section{Smallness with this same normalization}\label{sec:size}
The analytic gain comes from using a different polynomial basis to estimate
the determinant. Its first zeros remove the earliest atoms of the measure.
The signs of the character do not affect this argument, which bounds
absolute weights.

\begin{proposition}\label{prop:size}
There is a constant $C_0>0$, independent of $n$ and $m$, such that, whenever
\begin{equation}\label{eq:size-range}
 2m\leq(s-B^2-10)n,
\end{equation}
we have
\begin{equation}\label{eq:size-bound}
 H_{n,m}|\Delta_{n,m}|
 \leq (C_0m)^m A^{mn}e^{m^2+(2s+10B^2)mn}
       \left(\frac{10B^2 n}{m}\right)^{(s-B^2)mn-2m^2}.
\end{equation}
The constant may depend on the fixed parameters $s,B,A_0$ and $w_q$.
\end{proposition}
\begin{proof}
We give the estimate of \cite[Proposition~4.1]{FLSZ} with the character
weights included. Put $\nu_i=\lfloor i/(2h+2)\rfloor$ and take the monic
polynomials
\[
 V_i(z)=z^{i-(2h+2)\nu_i}
       \prod_{\theta\in\Theta\cup\{1\}}\prod_{v=0}^{\nu_i-1}
        \left(z^2-(v+\theta+n/2)^2\right).
\]
They have degree $i$, so $\Delta_{n,m}=\det(\Lcal_n(V_iV_j))$.
At every contributing point $t=v+j/q$ we have $\chi_4(j)\ne0$ and
$j/q\in\Theta$. Thus both $V_i(t+n/2)$ and $V_i(-t-n/2)$ vanish if
$v<\nu_i$. For $v\geq\nu_i$, all their roots in absolute value
are no larger than $t+n/2$, which gives
\[
 |V_i(t+n/2)|=|V_i(-t-n/2)|\leq(t+n)^i.
\]
Let $S_i$ be the sum defining $\Lcal_n(V_i^2)$ with every weight
$w_q\chi_4(j)$ replaced by its absolute value. Cauchy--Schwarz and the
determinant expansion give
\begin{equation}\label{eq:det-majorant}
 |\Lcal_n(V_iV_j)|\leq\sqrt{S_iS_j},\qquad
 |\Delta_{n,m}|\leq m!\prod_{i<m}S_i\leq m^m\prod_{i<m}S_i.
\end{equation}

Write $B_0=1+\max\mathcal Q_B$. All contributing points satisfy $t\geq1/(B_0-1)$.
The inequalities
\[
 (t)_{n+1}\geq\frac{n!}{B_0}(1+t/n)^{n+1},\qquad
 (t+\theta)_n\leq(t+n)^n,\qquad n!\geq(n/e)^n
\]
imply
\[
 R_n(t)\leq B_0^s A^n e^{hn}
             \left(\frac{n}{t+n}\right)^{(n+1)s-nh}.
\]
For $e_i=(n+1)s-nh-2i$, condition \eqref{eq:size-range} gives
$e_i-1>10n$. An integral comparison therefore yields
\[
 \sum_{v\geq\nu_i}\left(\frac{n}{v+j/q+n}\right)^{e_i}
 \leq\frac{2}{(1+\nu_i/n)^{e_i-1}}.
\]
Since $h+1\leq B^2$ and $n\geq2$, we also have
$1+\nu_i/n>(i+1)/(2B^2n)$. Set $E=(s-B^2)n-2m>0$.
When $2B^2n/(i+1)<1$ the preceding inequalities allow the exponent to
be reduced to $E$; when that ratio is at least one, the same bound follows
by first using $(1+\nu_i/n)^{-(e_i-1)}\leq1$. Thus
\[
 S_i\leq C_0 A^n e^{B^2n}n^{2i}
                 \left(\frac{2B^2n}{i+1}\right)^E,
 \qquad C_0=4A_0B_0^s\sum_{q\in\mathcal Q_B}q|w_q|.
\]
We used only $|\chi_4(j)|\leq1$. Combining with
\eqref{eq:det-majorant} and $m!\geq(m/e)^m$ gives
\begin{equation}\label{eq:raw-det-bound}
 |\Delta_{n,m}|\leq(C_0m)^m A^{mn}e^{B^2mn}n^{m(m-1)}
                \left(\frac{2eB^2n}{m}\right)^{Em}.
\end{equation}
Finally, the elementary bounds $d_n\leq e^{2n}$ and
\[
 \log E_{n,i}\geq i\log n-i-n
\]
give
\[
 H_{n,m}\leq n^{-m(m-1)}e^{m^2+2(s+1)mn}.
\]
Multiply this by \eqref{eq:raw-det-bound}, use $2e<10$ and
$B^2+2\leq10B^2$, and obtain \eqref{eq:size-bound}.
\end{proof}

\section{Nonvanishing for the signed character measure}\label{sec:nonvanishing}
The weights in \eqref{eq:functional} can change sign, so positivity of
individual entries or of their determinant is unavailable. Instead we use
the signed-measure criterion of \cite[Proposition~5.4]{FLSZ}. We include
its argument to expose exactly which nonvanishing hypothesis the character
twist must satisfy.

\subsection{A signed-measure criterion}
Let $X\subset\R$ be closed and discrete, and let $\mu=\sum_{x\in X}a_x\delta_x$
be a real signed measure. Suppose, for some positive integer $D$, that
\begin{equation}\label{eq:growth}
 \sum_{x\in X}|a_x|(1+|x|)^{D-1}<\infty.
\end{equation}
Extend $a_x$ by zero outside $X$. Its Cauchy transform and moments are
\[
 f(z)=\sum_{x\in X}\frac{a_x}{z-x},\qquad
 \mu_k=\sum_{x\in X}a_x x^k\quad(0\leq k<D).
\]
The series for $f$ is locally uniformly convergent off $X$, and its residue
at $x$ is $a_x$.

\begin{proposition}[Fischler--Lai--Sprang--Zudilin]\label{prop:signed}
Assume \eqref{eq:growth} and
\begin{equation}\label{eq:weight-recurrence}
 \mathfrak A(x)a_{x+1}=\mathfrak B(x)a_x\quad(x\in\R),
\end{equation}
where $\mathfrak A,\mathfrak B\in\C[z]$ and
$1\leq d=\max(\deg\mathfrak A,\deg\mathfrak B)\leq D$. Put
\begin{equation}\label{eq:C-polynomial}
 \mathfrak C(z)=\sum_{x\in X}a_x
 \left(\frac{\mathfrak A(z)-\mathfrak A(x-1)}{z-(x-1)}
       -\frac{\mathfrak B(z)-\mathfrak B(x)}{z-x}\right).
\end{equation}
If
\begin{equation}\label{eq:rational-equation}
 \mathfrak A(z)r(z+1)-\mathfrak B(z)r(z)=\mathfrak C(z)
\end{equation}
has no solution $r\in\C(z)$, then for every integer $M\geq1$ with
$2(M+d-1)-2<D$ there is an $m\in[M,M+d-1]$ such that
$\det(\mu_{i+j})_{0\leq i,j<m}\ne0$.
\end{proposition}
\begin{proof}
All summands in \eqref{eq:C-polynomial} are polynomials of degree at most
$d-1$, and \eqref{eq:growth} makes their coefficients absolutely summable.
Subtract the polynomial divided differences from
$\mathfrak A(z)f(z+1)-\mathfrak B(z)f(z)$. The remaining partial fractions
cancel by \eqref{eq:weight-recurrence}. Consequently
\begin{equation}\label{eq:f-equation}
 \mathfrak A(z)f(z+1)-\mathfrak B(z)f(z)=\mathfrak C(z).
\end{equation}

Set $N=M+d-1$. We first show that the $N\times M$ Hankel matrix
$(\mu_{i+j})_{0\leq i<N,\,0\leq j<M}$ has rank $M$. Otherwise there
is a nonzero polynomial $Q$ of degree at most $M-1$ such that
\[
 \sum_{x\in X}a_x Q(x)x^i=0\qquad(0\leq i<N).
\]
The polynomial
\[
 P(z)=\sum_{x\in X}a_x\frac{Q(z)-Q(x)}{z-x}
\]
has degree at most $M-2$ (and is zero if $M=1$). Expanding $1/(z-x)$
to order $N$ and using the vanishing moments gives
\[
 Q(z)f(z)-P(z)=\frac1{z^N}\sum_{x\in X}
                     a_x Q(x)x^{N-1}\frac{x}{z-x}.
\]
The absolute sum of $a_xQ(x)x^{N-1}$ is finite, since
$\deg Q+N-1\leq2N-2<D$. On each line $z=iy$ and $z=1+iy$ with
$y\geq1$, the last quotient is bounded in absolute value by $2$ and
tends to zero for each fixed $x$. Dominated convergence yields
\begin{equation}\label{eq:good-rational}
 Q(z)f(z)-P(z)=o(y^{-N})\qquad(y\to\infty)
\end{equation}
on both lines. The polynomial
\[
 \mathfrak A(z)P(z+1)Q(z)-\mathfrak B(z)P(z)Q(z+1)
       -\mathfrak C(z)Q(z)Q(z+1)
\]
is therefore $o(1)$ on the imaginary axis, by \eqref{eq:f-equation},
\eqref{eq:good-rational} and $N\geq d+\deg Q$. It vanishes identically.
Dividing by $Q(z)Q(z+1)$ makes $P/Q$ a solution of
\eqref{eq:rational-equation}, a contradiction.

It remains to pass from this rectangular rank to a principal determinant.
Here is the elementary Hankel fact used in \cite[Lemma~5.3]{FLSZ}. In
the $N\times N$ Hankel matrix, let $m$ be the largest number of linearly
independent initial columns. We have $m\geq M$. If $m=N$, we are done.
Otherwise the next column supplies a recurrence
\[
 \mu_{i+m}+\sum_{j=0}^{m-1}c_j\mu_{i+j}=0\qquad(0\leq i<N).
\]
For the row vectors of the first $m$ columns, this recurrence expresses
each row from index $m$ onward as a combination of the preceding $m$
rows. The indices needed in each coordinate lie within the displayed
recurrence. Thus the first $m$ rows span all these rows. Since their
column rank is $m$, their first $m\times m$ block is invertible. This
is the required principal determinant.
\end{proof}

\subsection{The two tails of the character measure}
Let $X=\Z+(\Theta\cup\{1\})$. On $X$, define the combined grid weight
\[
 w(x)=\sum_{\substack{q\in\mathcal Q_B\\qx\in\Z}}w_q\chi_4(qx),
 \qquad \widehat R_n(x)=R_n(x-n/2).
\]
Because each $q$ is divisible by four,
\begin{equation}\label{eq:w-parity}
 w(x+1)=w(x),\qquad w(-x)=-w(x),\qquad w(k)=0\ (k\in\Z).
\end{equation}
Equation~\eqref{eq:reflection} says that $\widehat R_n$ is even.
The moments in \eqref{eq:functional} are those of the signed measure with
masses
\begin{equation}\label{eq:mass}
 a_x=
 \begin{cases}
  A_0\sgn(x)w(x)\widehat R_n(x),&x\in X,\ |x|>n/2,\\
  0,&\text{otherwise}.
 \end{cases}
\end{equation}
On the positive tail this follows by putting $x=v+j/q+n/2$; $n/2$ is
an integer. On the negative tail, the oddness of $w$ and the extra factor
$\sgn(x)$ make the reflected mass equal to the positive one, just as the
two evaluations of $P$ in \eqref{eq:functional} require. The sign factor
cannot be dropped. The measure is even, but its individual masses can
have either sign.

The ratio of consecutive kernel values is
\begin{equation}\label{eq:kernel-ratio}
 \frac{\widehat R_n(z+1)}{\widehat R_n(z)}
   =\frac{\mathcal B_n(z)}{\mathcal A_n(z)},
\end{equation}
where
\begin{align}
 \mathcal A_n(z)&=(z+1+n/2)^{s+1}(z-n/2)
                    \prod_{\theta\in\Theta}(z+\theta-n/2),\label{eq:A-rec}\\
 \mathcal B_n(z)&=(z-n/2)^{s+1}(z+1+n/2)
                    \prod_{\theta\in\Theta}(z+\theta+n/2).\label{eq:B-rec}
\end{align}
These are the coefficients of \cite[(5.11)]{FLSZ} with the products
taken over our $\Theta$. Both have degree
\[
 d=s+h+2.
\]
Away from the central gap, \eqref{eq:kernel-ratio} and periodicity give
$\mathcal A_n(x)a_{x+1}=\mathcal B_n(x)a_x$.
At a transition into the positive tail, write
$x=n/2-1+\theta$, where $0<\theta<1$ is a grid fraction. The factor
indexed by $1-\theta$ in $\mathcal A_n(x)$ vanishes. At a transition out
of the negative tail, $x=-n/2-\theta$, and the factor indexed by $\theta$
in $\mathcal B_n(x)$ vanishes. Integral endpoints carry no mass by
\eqref{eq:w-parity}. Points outside $X$ remain outside $X$ after an
integer shift. Thus the recurrence holds for the extension by zero at
every real $x$.

Since $\deg R_n=nh-(n+1)s$, condition \eqref{eq:growth} holds with
\begin{equation}\label{eq:moment-D}
 D=(n+1)s-nh-1.
\end{equation}
Indeed the summands at infinity are $O(|x|^{-2})$, and $X$ is a finite
union of translates of $\Z$.

\begin{lemma}\label{lem:no-rational}
For this measure and the polynomials \eqref{eq:A-rec}--\eqref{eq:B-rec},
the rational equation in Proposition~\ref{prop:signed} has no solution.
\end{lemma}
\begin{proof}
Suppose $r\in\C(z)$ solves it. The Cauchy transform $f$ solves the same
equation. Subtracting and using \eqref{eq:kernel-ratio}, we find that
\[
 g(z)=\frac{f(z)-r(z)}{\widehat R_n(z)}
 \quad\text{satisfies}\quad g(z+1)=g(z)
\]
as a meromorphic identity.

Choose the largest denominator with a nonzero elimination weight,
$q_* =\max\{q:w_q\ne0\}$, and set $\theta_*=1/q_*$. At either point
$x=\pm T+\theta_*$, with integral $T$, a grid denominator $q$ contributes
only if $q_*\mid q$. Maximality leaves just $q=q_*$ among the nonzero
weights. Moreover $\chi_4(q_*(\pm T)+1)=1$, so
\[
 w(T+\theta_*)=w(-T+\theta_*)=w_{q_*}\ne0.
\]
Reflection pairs $T+\theta_*$ with $-T-\theta_*$, where the masses
agree. Periodicity pairs $T+\theta_*$ with $-T+\theta_*$, where
the grid weight is the same and the sign of $x$ changes.
Choose $T$ large enough that these latter points lie outside the central gap
and avoid the finitely many poles of $r/\widehat R_n$ and all zeros
or poles of $\widehat R_n$. From \eqref{eq:mass}, the residues of $g$
at these two points are respectively
\[
 A_0w_{q_*}\quad\text{and}\quad -A_0w_{q_*}.
\]
But the points differ by the integer $2T$, so a one-periodic meromorphic
function must have equal residues there. This is a contradiction.
\end{proof}

\begin{corollary}\label{cor:nonzero}
Let $M\in\Z_{\geq1}$. If $d\leq D$ and $2(M+d-1)-2<D$, then some integer
$m\in[M,M+d-1]$ satisfies $\Delta_{n,m}\ne0$.
\end{corollary}
\begin{proof}
Apply Proposition~\ref{prop:signed} to \eqref{eq:mass}, using
Lemma~\ref{lem:no-rational}. Its moments are exactly those defining
$\Delta_{n,m}$.
\end{proof}
\section{The asymptotic comparison}\label{sec:finish}
We now choose the parameters. The number of usable grids grows linearly
with $B$, while $\log A(B)$ grows like $B^2\log B$. The determinant
estimate permits $B$ to be larger by a factor $\sqrt{\log\log s}$ than
the earlier single-form scale. The following calculation makes the gain
and the order of limits explicit.

Fix
\[
 0<\lambda<\sqrt{\frac{12}{\kappa}},\qquad
 B=\lambda Y(s),\qquad C=\frac{s}{\log\log s}.
\]
For all sufficiently large even $s$, the hypotheses on $B$ and $s$ in
Sections~\ref{sec:functional}--\ref{sec:size} hold, and
\[
 2C<s-B^2-10.
\]
Fix one such $s$. As $n$ tends to infinity through even integers, take
$M=\lceil Cn\rceil$. The degree conditions in Corollary~\ref{cor:nonzero}
hold for all sufficiently large $n$, since $h\leq B^2-1$ and
$2C<s-B^2-10$. Consequently we may choose
\[
 M\leq m=m(n)\leq M+d-1,
 \qquad\Delta_{n,m}\ne0,
 \qquad\frac{m}{n}\longrightarrow C.
\]
The same strict inequality ensures \eqref{eq:size-range} for these $m$,
once $n$ is large. Proposition~\ref{prop:integerizer} then gives nonzero
integers
\[
 J_n=H_{n,m(n)}\Delta_{n,m(n)},\qquad |J_n|\geq1.
\]
We will contradict this using Proposition~\ref{prop:size}.

Take logarithms in \eqref{eq:size-bound}, divide by $mn$, and first let
$n\to\infty$ with $s,B,C$ fixed. The term $\log(C_0m)/n$ disappears,
regardless of the size of the fixed elimination weights or $A_0$. We obtain
\begin{equation}\label{eq:rate}
 \limsup_{\substack{n\to\infty\\2\mid n}}
 \frac{\log|J_n|}{m(n)n}
 \leq \log A+C+2s+10B^2
        +(s-B^2-2C)\log\frac{10B^2}{C}.
\end{equation}
The right-hand side is now an expression in $s$. As $s\to\infty$,
\eqref{eq:A1}--\eqref{eq:A2} and $\log B\sim\frac12\log s$ give
\begin{equation}\label{eq:arithmetic-rate}
 \log A\leq
 \left(\frac{\kappa\lambda^2}{12}+o(1)\right)s\log\log s.
\end{equation}
For example, the $A_2$ contribution is
$O(s(\log\log s)^3/\log s)=o(s\log\log s)$, so it does not affect
the leading constant. On the analytic side,
\begin{align*}
 \log\frac{10B^2}{C}
 &=-\log\log s+2\log\log\log s+\log(10\lambda^2),\\
 s-B^2-2C&=s(1-o(1)).
\end{align*}
Inserting these into \eqref{eq:rate} gives
\[
 \limsup_{\substack{n\to\infty\\2\mid n}}
 \frac{\log|J_n|}{m(n)n}
 \leq\left(\frac{\kappa\lambda^2}{12}-1+o(1)\right)s\log\log s.
\]
The coefficient is negative for every sufficiently large $s$, by the
choice of $\lambda$. For each such fixed $s$, we therefore obtain
$|J_n|<1$ for all sufficiently large even $n$, contradicting
$J_n\in\Z\setminus\{0\}$.

The contradiction disproves the assumption $I(s)<|\mathcal Q_B|$.
Lemma~\ref{lem:grid-count} now implies
\[
 \liminf_{\substack{s\to\infty\\2\mid s}}\frac{I(s)}{Y(s)}
 \geq\frac{\lambda\kappa}{3}.
\]
Letting $\lambda$ increase to $\sqrt{12/\kappa}$ gives
$(\kappa/3)\sqrt{12/\kappa}=2\sqrt{\kappa/3}$, proving Theorem~\ref{thm:main}.

\begin{remark}
The proof concerns growing intervals and uses no explicit threshold for
$s$. In particular, it does not by itself give an irrationality statement
for a short prescribed list or for an individual even beta value. Both
the character-dependent nonvanishing argument and the fixed-$s$ order of
limits are essential to the conclusion above.
\end{remark}

\paragraph{Verification.}
Independent AI models audited the mathematics and exposition of the original and
revised proofs. The first proposed the restricted grid.

\begin{thebibliography}{9}

\bibitem{FLSZ}
S.~Fischler, L.~Lai, J.~Sprang and W.~Zudilin,
\emph{Many more odd zeta values are irrational},
arXiv:2611.16156v1, 12 November 2026.
\url{https://arxiv.org/abs/2611.16156v1}.

\bibitem{KT}
V.~Kumar and R.~Tosi,
\emph{Irrationality of values of $L$-functions for real Dirichlet characters},
arXiv:2610.41458v1, 30 October 2026.
\url{https://arxiv.org/abs/2610.41458v1}.

\bibitem{LY}
L.~Lai and P.~Yu,
\emph{A note on the number of irrational odd zeta values},
Compositio Mathematica \textbf{156} (2020), no.~8, 1699--1717.
\url{https://arxiv.org/abs/1911.08458v2}.

\bibitem{THcount}
twinheterodox,
\emph{Explicit lower bounds for the number of irrational even beta values},
Hexagon 2611.00139v1, 4 November 2026.
\url{https://hexagonmath.org/2611.00139v1}.

\bibitem{beta4repo}
twinheterodox, \emph{Irrationality and linear independence of Dirichlet beta values},
GitHub repository, 2026. \url{https://github.com/twinheterodox/beta4}.

\end{thebibliography}
\end{document}
